\section{Related Work}
\label{sec:related}

\textbf{Market efficiency and Bitcoin.} The efficient market hypothesis \cite{fama1970} implies that publicly available price information cannot be used to forecast returns reliably. For Bitcoin, early evidence pointed to inefficiency \cite{urquhart2016}, but simple transformations of returns already restored weak-form efficiency in follow-up tests \cite{nadarajah2017}, and Urquhart himself reported that the market was becoming more efficient over time. This tension makes Bitcoin a natural test bed: if machine learning has an exploitable edge anywhere, a young, retail-heavy, 24/7 market is a plausible candidate.

\textbf{Machine learning for Bitcoin forecasting.} A large literature applies LSTMs, gradient boosting and other learners to Bitcoin prices, often reporting high accuracy (e.g., \cite{mcnally2018}). Jaquart et al.~\cite{jaquart2021} compare several model classes on short horizons and find predictability that is statistically detectable but small and largely eroded by transaction costs. Many studies, however, evaluate on price levels, where predicting yesterday's price already yields low error; use random train/test splits that leak future information; report a single test window; or compare a proposed model only against weak or untuned alternatives. Kapoor and Narayanan \cite{kapoor2023} document how such data leakage has inflated published results across machine-learning-based science.

\textbf{Evaluation methodology.} We draw on established forecasting and finance practice. Out-of-sample $R^2$ against a naive benchmark \cite{campbell2008, welch2008} measures whether a model beats the random walk. The Diebold--Mariano test \cite{diebold1995} assesses whether accuracy differences are statistically significant. Purged and embargoed walk-forward validation \cite{lopezdeprado2018} prevents leakage through overlapping labels, and cross-validation for time series requires respecting temporal order \cite{bergmeir2012}. When many models or strategies are compared, multiple-testing corrections \cite{holm1979, harvey2016} and the deflated Sharpe ratio \cite{bailey2014} guard against declaring the luckiest backtest a discovery.

Our contribution is to combine these safeguards in a single, open, reproducible benchmark spanning 32 models from six families plus the random walk, evaluated on the same leakage-free folds over seven years that include several distinct market regimes.
